BioCLIP is a vision--language foundation model that names organisms across the whole tree of life, and it is used
to support biodiversity monitoring. We ask how robust it is to realistic images it has never seen. We build
\emph{BioFool}, a generative application that produces faithful images of a given species that BioCLIP's public
open-set classifier ($\sim$450k species) gets wrong. Fidelity is judged by GPT-4o, a model of another family, and
by humans. We adapt Stable Diffusion in four ways: context engineering of the prompt (including an LLM attacker),
editing of real photographs, inference-time search with a new \emph{early-rejection} scheme, and post-training of a
LoRA with reinforcement learning (DDPO with group-relative advantages). We then accelerate sampling. On four datasets
(synthetic images, iNaturalist photos, invasive species and BioCLIP's own Rare Species benchmark), \todo{headline
result: e.g. DDPO raises the attack success rate from X\% to Y\% on unseen species; early rejection keeps Z\% of
best-of-N's gain at 55\% of its cost}. Code: \url{\todo{repository URL}}.
